%\documentclass[a4paper, 14pt]{article}
\documentclass[a4paper, 14pt]{extarticle} 

\usepackage[utf8]{inputenc}
\usepackage[T2A]{fontenc}
\usepackage[english,russian]{babel}
% \fontsize{14}{17}

% \usepackage[none]{hyphenat}  % Полное отключение переносов

\usepackage{setspace}
%\полуторный интервал
% \onehalfspacing
\usepackage[left=30mm, top=20mm, right=15mm, bottom=20mm, nohead, nofoot]{geometry}
\usepackage{amsmath,amsfonts,amssymb, gensymb, tikz, graphicx}
\usepackage{fancybox,fancyhdr}
\usepackage{multicol}
\usepackage{indentfirst}
\usepackage{tabto}
\usepackage{titlesec}
\usepackage{listings}
\usepackage{fancyhdr}
\usepackage{booktabs}
\usepackage{float}
\usepackage{tocloft} % для настройки оглавления
\usepackage{titlesec} % для настройки заголовков
\usepackage{hyperref} % переносит длинные ссылки
\pagestyle{fancy}
\fancyhf{}  % Очистить стандартные колонтитулы
\renewcommand{\headrulewidth}{0pt} % Убираем линию под верхним колонтитулом
\renewcommand{\footrulewidth}{0pt} % Убираем линию над нижним колонтитулом (если нужно)
\cfoot{\thepage}  % Номер справа внизу
% \pagestyle{fancy}
% \fancyhf{}
% \fancyhead[R]{A. Khalitov}
% \fancyfoot[R]{\thepage}
% \fancyhead[L]{Срез знаний за 7 класс}
\usepackage[ruled,vlined]{algorithm2e}
\usepackage{algorithm, algorithmic, algpseudocode}
\setcounter{page}{4}
% \headsep=10mm
\usepackage{xcolor}
% \usepackage[style=gost-numeric]{biblatex} % для оформленния ссылок на литературу по ГОСТу
\usepackage{hyperref}
\hypersetup{colorlinks=true, allcolors=black}
\newcommand{\lr}[1]{\left( {#1} \right)}

% Включаем отступ после \subsection
\makeatletter
\patchcmd{\@startsection}{\@afterindentfalse}{\@afterindenttrue}{}{}
\makeatother

\usepackage{csquotes} % требуется biblatex
\usepackage[style=numeric,sorting=none]{biblatex}
\addbibresource{sources.bib}

% Настройка — оставить только название источника в списке литературы
\AtEveryBibitem{%
	\clearfield{author}%
	\clearfield{editor}%
	\clearfield{publisher}%
	\clearfield{year}%
	\clearfield{journal}%
	\clearfield{note}%
}

% Переопределим шаблоны, чтобы показывать только заголовок
\DeclareFieldFormat{title}{#1}
\DeclareBibliographyDriver{book}{\printfield{title}}
\DeclareBibliographyDriver{article}{\printfield{title}}
\DeclareBibliographyDriver{online}{\printfield{title}}

% \usepackage[backend=biber, style=numeric, sorting=none]{biblatex}

% % Настраиваем отображение записей в библиографии
% \DeclareFieldFormat[article,book]{title}{#1} % Название выводится первым
% \renewbibmacro{in:}{} % Убираем "In:" для статей

% % Настраиваем отображение ссылок
% \DeclareFieldFormat{labelnumberwidth}{[#1]} % Формат ссылок: [1], [2], ...

% Отступы
% Абзацный отступ
\setlength{\parindent}{1.25cm}
% Настройка отступов для section
\titleformat{\section}
{\normalfont\Large\bfseries}{\thesection}{1em}{} % {формат}{метка}{отступ от метки}{перед заголовком}[после заголовка]
\titlespacing*{\section}{1.25cm}{3.5ex}{2.3ex} % {левый отступ}{отступ сверху}{отступ снизу}

% Настройка отступов для subsection
\titleformat{\subsection}
{\normalfont\large\bfseries}{\thesubsection}{1em}{}
\titlespacing*{\subsection}{1.25cm}{2.5ex}{1.5ex}

% Настройка отступов для subsubsection
\titleformat{\subsubsection}
{\normalfont\large\bfseries}{\thesubsubsection}{1em}{}
\titlespacing*{\subsubsection}{1.25cm}{2.5ex}{1.5ex}
% \titleformat{\section}{\normalfont\normalsize\bfseries}{\thesection}{1em}{}
% \titleformat{\subsection}{\normalfont\normalsize\bfseries}{\thesubsection}{1em}{}

% Настройка содержания
\renewcommand{\cfttoctitlefont}{\bfseries\fontsize{14}{16}\MakeUppercase}
\renewcommand{\cftaftertoctitle}{\centering} % выравнивание заголовка

% Настройка отступов и стиля пунктов
\renewcommand{\cftsecindent}{0pt} % отступ для разделов
\renewcommand{\cftsubsecindent}{0cm} % отступ для подразделов
\renewcommand{\cftsubsubsecindent}{0cm} % отступ для пунктов

% Настройка точек и шрифта
\renewcommand{\cftsecfont}{\normalfont} % шрифт разделов
\renewcommand{\cftsubsecfont}{\normalfont} % шрифт подразделов
\renewcommand{\cftsubsubsecfont}{\normalfont} % шрифт пунктов

\renewcommand{\cftsecleader}{\cftdotfill{\cftdotsep}} % точки для разделов
\renewcommand{\cftsubsecleader}{\cftdotfill{\cftdotsep}} % точки для подразделов
\renewcommand{\cftsubsubsecleader}{\cftdotfill{\cftdotsep}} % точки для пунктов

% Общий интервал между всеми пунктами
\setlength{\cftbeforetoctitleskip}{0pt} % отступ перед заголовком "Содержание"
\setlength{\cftaftertoctitleskip}{0pt}  % отступ после заголовка "Содержание"
% Расстояние между разделами (section)
\setlength{\cftbeforesecskip}{0pt} % где X - нужное значение в пунктах
% Расстояние между подразделами (subsection)
\setlength{\cftbeforesubsecskip}{0pt}
% Расстояние между подподразделами (subsubsection)
\setlength{\cftbeforesubsubsecskip}{0pt}

% Настройка номеров страниц
\renewcommand{\cftsecpagefont}{\normalfont}   % для разделов (section)
\renewcommand{\cftsubsecpagefont}{\normalfont} % для подразделов (subsection)
\renewcommand{\cftsubsubsecpagefont}{\normalfont} % для пунктов (subsubsection)

% Для листинга кода
\lstset{
	language=C++,
	basicstyle=\small\ttfamily,
	keywordstyle=\color{blue}\bfseries,
	stringstyle=\color{red},
	commentstyle=\color{green!50!black}\textit,
	% numbers=left,
	% numberstyle=\tiny\color{gray},
	% stepnumber=1,
	% numbersep=5pt,
	showspaces=false,
	showstringspaces=false,
	% frame=single,
	tabsize=2,
	breaklines=true,
	backgroundcolor=\color{white},
}

\lstset{
	language=bash,
	basicstyle=\ttfamily\footnotesize,
	keywordstyle=\color{blue}\bfseries,
	commentstyle=\color{green!60!black}\itshape,
	stringstyle=\color{red!70!black},
	numberstyle=\tiny\color{gray},
	%numbers=left,
	numbersep=5pt,
	stepnumber=1,
	backgroundcolor=\color{gray!8},
	showspaces=false,
	showstringspaces=false,
	showtabs=false,
	tabsize=4,
	breaklines=true,
	breakatwhitespace=true,
	frame=single,
	framerule=0.5pt,
	framesep=6pt,
	captionpos=b,
	morekeywords={sudo, chmod, chown, apt, yum, brew, git, docker, systemctl, 
		export, source, alias, unset, trap, set, shopt},
	sensitive=false,
}

\usepackage{pgfplots}
\pgfplotsset{compat=1.18}

\usepackage{array}
\newcolumntype{C}[1]{>{\centering\arraybackslash}p{#1}}

% \bibliographystyle{gost-numeric.bbx}
% \usepackage[parentracker=true,
% backend=biber,
% hyperref=false,
% bibencoding=utf8,
% style=numeric-comp,
% language=auto,
% autolang=other,
% citestyle=gost-numeric,
% defernumbers=true,
% bibstyle=gost-numeric,
% sorting=ntvy,
% ]{biblatex}
% \addbibresource{library.bib}

% \documentclass[a4paper,12pt]{article}
% \usepackage[utf8]{inputenc}
% \usepackage[russian]{babel}
% \usepackage{amsmath,amssymb}
% \usepackage{times}
% \usepackage{setspace}
% \usepackage{graphicx}
% \usepackage{float}
% \fontsize{14}{17}
% \usepackage{geometry}
% \usepackage{etoolbox} % Для модификации команд
% \setlength{\parindent}{2cm} % Устанавливаем отступ в преамбуле

% \geometry{
	%   left=1.5cm,
	%   right=1.5cm,
	%   top=2cm,
	%   bottom=2cm
	% }

% % Включаем отступ после \subsection
% \makeatletter
% \patchcmd{\@startsection}{\@afterindentfalse}{\@afterindenttrue}{}{}
% \makeatother


\begin{document}
	
	\begin{titlepage}
		\centering
		
		\textbf{Министерство науки и высшего образования РФ} \\
		\textbf{Федеральное государственное бюджетное образовательное учреждение высшего образования} \\
		\textbf{«Уфимский университет науки и технологий»} \\[0.5cm]
		
		\textbf{Кафедра} Высокопроизводительных вычислений и дифференциальных \\[0.2cm] уравнений \\[0.1cm]
		
		\begin{table}[H]
			\noindent
			\hspace{5.5cm}
			\footnotesize
			\renewcommand{\arraystretch}{1.49}
			\begin{tabular}{|p{0.5cm}|p{0.4cm}|p{0.4cm}|p{0.4cm}|p{0.4cm}|p{0.4cm}|p{0.4cm}|p{0.4cm}|p{0.4cm}|p{0.4cm}|p{0.4cm}|}
				\hline
				& 1 & 2 & 3 & 4 & 5 & 6 & 7 & 8 & 9 & 10 \\ \hline
				100 & & & & & & & & & & \\ \hline
				90 & & & & & & & & & & \\ \hline
				80 & & & & & & & & & & \\ \hline
				70 & & & & & & & & & & \\ \hline
				60 & & & & & & & & & & \\ \hline
				50 & & & & & & & & & & \\ \hline
				40 & & & & & & & & & & \\ \hline
				30 & & & & & & & & & & \\ \hline
				20 & & & & & & & & & & \\ \hline
				10 & & & & & & & & & & \\ \hline
				0 & & & & & & & & & & \\ \hline
			\end{tabular}
		\end{table}
		
		\vspace{-0.2cm}
		
		{\textbf{ОТЧЕТ}} \\[0.6cm]
		{\textbf{По лабораторной работе №1} \\[0.3cm]
			{Параллельное вычисление суммы числового ряда} \\[0.3cm]
			{По дисциплине \textbf{...}}}
		
		\vfill
		\vspace{1cm}
		\noindent
		\begin{table}[H]
			\hspace{-0.8cm}
			\renewcommand{\arraystretch}{1.5} 
			\begin{tabular}{|C{3.1cm}|C{4cm}|C{2.4cm}|C{2.2cm}|C{3.2cm}|}
				\hline
				Группа & Фамилия И.О. & Подпись & Дата & Оценка \\ \hline
				ПМ-401Б & & & & \\ \hline
				Студент & Халитов А.Ю. & & & --- \\ \hline
				Принял &  & & & \\ \hline
			\end{tabular}
		\end{table}
		\vspace{0.25cm}
		Уфа 2026
	\end{titlepage}
	
	\section{Цель работы}
	Для многопроцессорных вычислительных систем с распределенной памятью на примере задачи параллельного вычисления суммы числового ряда научиться программно реализовывать простейшие параллельные вычислительные алгоритмы с использованием базовых функций библиотеки MPI, а также проводить анализ их эффективности.
	
	\section{Постановка задачи}
	Требуется вычислить сумму числового ряда:
	\[
	S = \sum_{n=1}^{N} a_n = a_1 + a_2 + \dots + a_N, \quad a_n \in \mathbb{R}, \; n \in \mathbb{N}
	\]
	В распоряжении находится многопроцессорная вычислительная система (МВС) с $P$ процессорами (ядрами).
	
	\section{Алгоритм параллельного вычисления}
	Для реализации параллельного алгоритма используются базовые функции MPI:
	\texttt{MPI\_Init}, \texttt{MPI\_Finalize}, \texttt{MPI\_Comm\_size}, \texttt{MPI\_Comm\_rank}, \texttt{MPI\_Send}, \texttt{MPI\_Recv}, \texttt{MPI\_Wtime}, \texttt{MPI\_Barrier}.
	
	Схема работы параллельной программы:
	\begin{enumerate}
		\item Ввод-вывод данных осуществляется через нулевой процесс ($\texttt{Rank} = 0$), который передает остальным процессам число членов ряда $N$, полученное из аргументов командной строки.
		\item Время начала выполнения параллельного участка программы замеряется с помощью функции \texttt{MPI\_Wtime()}.
		\item Каждый процесс определяет количество элементов для суммирования:
		\[
		h = \lfloor N / P \rfloor, \qquad m = N \bmod P
		\]
		Если остаток $m \neq 0$, остаток равномерно распределяется между процессорами.
		\item Каждый процесс считывает и рассчитывает выделенную ему часть ряда.
		\item Процессы с номерами $1 \dots P-1$ пересылают свои частичные суммы нулевому процессу с помощью блокирующей функции \texttt{MPI\_Send}.
		\item Нулевой процесс принимает данные с помощью \texttt{MPI\_Recv}, суммирует их со своей частичной суммой, фиксирует время окончания параллельного участка и выводит итоговую сумму и затраченное время $T_p$.
	\end{enumerate}
	
	\section{Результаты экспериментов}
	Эксперименты проводились на суперкомпьютерном кластере при числе процессов $P \in \{1, 2, 4, 8, 16\}$. Число членов ряда $N = \text{[значение]}$ подбиралось таким образом, чтобы время выполнения на одном процессоре ($P=1$) составляло порядка 30 cекунд.
	
	\begin{table}[H]
		\centering
		\begin{tabular}{|c|c|}
			\hline
			$P$ & $T_p$ \\ \hline
			1   & 30  \\ \hline
			2   & --- \\ \hline
			4   & --- \\ \hline
			8   & --- \\ \hline
			16  & --- \\ \hline
			32  & --- \\ \hline
			64  & --- \\ \hline
			128 & --- \\ \hline
		\end{tabular}
		\caption{Время выполнения параллельного участка кода $T_p$ в секундах}
		\label{tab:time}
	\end{table}
	
	\section{Анализ ускорения и эффективности}
	Ускорение $S_p^*$ и эффективность $E_p^*$ рассчитываются по формулам:
	\[
	S_p^* = \frac{T_1^*}{T_p}, \qquad E_p^* = \frac{S_p^*}{P}
	\]
	где $T_1^*$ — время работы программы на одном процессоре, $T_p$ — время работы программы на $P$ процессорах.
	
	\begin{table}[H]
		\centering
		\begin{tabular}{|c|c|c|}
			\hline
			$P$ & {Ускорение $S_p^*$} & Эффективность $E_p^*$ \\ \hline
			1   & 1.00 & 1.00 \\ \hline
			2   & ---  & ---  \\ \hline
			4   & ---  & ---  \\ \hline
			8   & ---  & ---  \\ \hline
			16  & ---  & ---  \\ \hline
			32  & ---  & ---  \\ \hline
			64  & ---  & ---  \\ \hline
			128 & ---  & ---  \\ \hline
		\end{tabular}
		\caption{Расчетные значения ускорения $S_p^*$ и эффективности $E_p^*$}
		\label{tab:speedup_eff}
	\end{table}
	
	\subsection{Графики зависимостей ускорения и эффективности}
	
	На рисунке~\ref{fig:graphs} приведена зависимость ускорения и эффективности от числа процессоров $P$.
	
	\begin{figure}[H]
		\centering
		\begin{minipage}{0.48\textwidth}
			\centering
			\begin{tikzpicture}
				\begin{axis}[
					xlabel={Число процессов $P$},
					ylabel={Ускорение $S_p^*$},
					grid=major,
					xmode=log,
					log basis x={2},
					xmin=1, xmax=128,
					legend pos=north west
					]
					% График для N1
					\addplot[mark=*, blue] coordinates {
						(1,1) (2,1.95) (4,3.8) (8,7.2) (16,13.5) (32,24.0) (64,40.0) (128,60.0)
					};
					\addlegendentry{$N_1$}
					
					% График для N2
					\addplot[mark=square, red] coordinates {
						(1,1) (2,1.98) (4,3.9) (8,7.5) (16,14.2) (32,26.5) (64,48.0) (128,75.0)
					};
					\addlegendentry{$N_2$}
				\end{axis}
			\end{tikzpicture}
		\end{minipage}
		\hfill
		\begin{minipage}{0.48\textwidth}
			\centering
			\begin{tikzpicture}
				\begin{axis}[
					xlabel={Число процессов $P$},
					ylabel={Эффективность $E_p^*$},
					grid=major,
					xmode=log,
					log basis x={2},
					xmin=1, xmax=128,
					ymin=0, ymax=1.1,
					legend pos=south west
					]
					% График для N1
					\addplot[mark=*, blue] coordinates {
						(1,1.0) (2,0.97) (4,0.95) (8,0.90) (16,0.84) (32,0.75) (64,0.62) (128,0.46)
					};
					\addlegendentry{$N_1$}
					
					% График для N2
					\addplot[mark=square, red] coordinates {
						(1,1.0) (2,0.99) (4,0.97) (8,0.93) (16,0.88) (32,0.82) (64,0.75) (128,0.58)
					};
					\addlegendentry{$N_2$}
				\end{axis}
			\end{tikzpicture}
		\end{minipage}
		\caption{Зависимость ускорения $S_p^*$ и эффективности $E_p^*$ от числа процессов $P$}
		\label{fig:graphs}
	\end{figure}
	
	\section*{Выводы}
	В ходе выполнения лабораторной работы:
	\begin{enumerate}
		\item Изучены базовые функции библиотеки MPI (\texttt{MPI\_Init}, \texttt{MPI\_Send}, \texttt{MPI\_Recv}, \texttt{MPI\_Wtime}) для взаимодействия процессов в системах с распределенной памятью.
		\item Разработана параллельная программа на C для вычисления суммы числового ряда.
		\item Проведены эксперименты при заданном числе процессов $P \in \{ 1, 2, \ldots \}$.
		\item Проанализированы показатели ускорения и эффективности, видно, что ... 
		
		установлено, что увеличение размера задачи приводит к росту эффективности параллельной обработки при большом числе процессоров.
	\end{enumerate}
	
	\newpage
	\section*{Листинг программы}
	
\end{document}